\documentclass[12pt]{report}
\usepackage[a4paper, margin=2.5cm]{geometry}
\pagenumbering{gobble}
\usepackage{hyperref}
\usepackage{tcolorbox}
\newcommand{\student}[3]{\begin{tcolorbox}\noindent Introduction to
    Psycholinguistics (WS20) \\ Report of \textbf{Paper~#3} by
    \textbf{#1 (#2)} 
  \end{tcolorbox}}



\begin{document}

%% Change:
\student{Shawon Ashraf}{3460184}{9}


\paragraph{APA reference:} Van Petten, C., & Rheinfelder, H. (1995). Conceptual relationships between spoken words and environmental sounds: Event-related brain potential measures. \textit{Neuropsychologia}, 33(4), 485-508.

\paragraph{Research Question(s):}
\begin{enumerate}
    \item Can words paired with related sound result in faster lexical decision making?
    \item Does context affect finding meaning of words when paired with environmental sound?
\end{enumerate}
\vfill
\paragraph{Methodology:}
There are two experiments in the whole process. The first one aims to answer the question if related sound affects in the way subjects can make lexical decisions about words. The second experiment explores the effects of context and environmental sound on the same decision process. The materials in use primarily consisted of non-speech sounds, audio recordings from commercial tapes, and words or non-words associated with non-speech sounds. \\
\noindent
Subjects participating in the first experiment were asked to listen to a stimulus created from the materials and then decide whether the stimulus presented next was a word or a non-word. The second experiment involved recording EEG of the subjects' to measure the ERP while they make a lexical decision.
\vfill

\paragraph{Results:}
The results from the first experiment show that accuracy for the lexical decision task were higher when words were paired with related sounds than when they were not (99.5\% vs 94.3\%). The EEG graphs generated in the second experiment show significant effects of context in making lexical decision for words when related environmental sound is paired. Together, the results indicate strong relation between conceptual relationship of words with related sounds can influence meaning inference.
\vfill

\paragraph{Personal Opinion:} 
The concepts of EEG and ERP are new to me and I found the procedures interesting, especially to measure overt attention and how data from these procedures can be used in measuring the performance of subjects in a lexical decision making task. However, the authors do not clearly summarize their results to make a combined formulation in the discussion part. Although they provide a detailed analysis of the results, especially ERP from EEG, it does not lead to a clear concise point.
\vfill


\end{document}

